% ----------------------------------------------------------
% Capítulo 3, Desenvolvimento da Solução (arquivo cap04-aprimoramento.tex)
% ----------------------------------------------------------
\chapter{Desenvolvimento da Solução}
\label{cap:aprimoramento}

O projeto seguiu a Estrutura Analítica do Projeto (EAP) da Figura~\ref{fig:eap}, que cobre a construção do sistema e a sua avaliação. Este capítulo descreve a construção, e o Capítulo~\ref{cap:metodologia} detalha a avaliação.

\begin{figure}[h]
\centering
\begin{tikzpicture}[
  font=\footnotesize,
  node distance=0.5cm,
  eapstep/.style={rectangle, draw, rounded corners, minimum width=13cm, minimum height=0.9cm, align=center, fill=blue!8},
  sub/.style={rectangle, draw, rounded corners, minimum width=2.4cm, minimum height=1.0cm, align=center, fill=gray!6, font=\scriptsize, inner sep=2pt},
  arr/.style={-{Stealth[length=2.2mm]}, thick}
]
  \node[eapstep] (a) {\textbf{1. Análise do Protótipo} \;\;\small(arquitetura, dependências, limitações documentadas)};
  \node[eapstep, below=of a] (r) {\textbf{2. Elicitação de Requisitos} \;\;\small(funcionais e não funcionais)};
  \node[eapstep, below=of r] (p) {\textbf{3. Planejamento da Migração} \;\;\small(limitação $\to$ solução; reaproveitado / reescrito / novo)};
  \node[eapstep, below=of p] (d) {\textbf{4. Desenvolvimento e Avaliação}};
  \node[sub, below=0.9cm of d, xshift=-5.2cm] (s1) {API\\(FastAPI)};
  \node[sub, right=0.2cm of s1] (s2) {Orquestração\\(LangChain)};
  \node[sub, right=0.2cm of s2] (s3) {Ferramentas\\de Busca};
  \node[sub, right=0.2cm of s3] (s4) {Persistência\\(PostGIS)};
  \node[sub, right=0.2cm of s4] (s5) {\textit{Pipeline}\\de Avaliação\\e Experimentos};
  \node[eapstep, below=1.3cm of s3] (c) {\textbf{Cronograma} \;\;\small(abr.--out./2026, fases F0--F7)};
  \draw[arr] (a) to (r);
  \draw[arr] (r) to (p);
  \draw[arr] (p) to (d);
  \draw[arr] (d.south) to ++(0,-0.30) -| (s1.north);
  \draw[arr] (d.south) to ++(0,-0.30) -| (s2.north);
  \draw[arr] (d.south) to ++(0,-0.30) -| (s3.north);
  \draw[arr] (d.south) to ++(0,-0.30) -| (s4.north);
  \draw[arr] (d.south) to ++(0,-0.30) -| (s5.north);
  \draw[arr] (d.east) to ++(0.8,0) |- (c.east);
\end{tikzpicture}
\caption{Estrutura Analítica do Projeto (EAP)}
\label{fig:eap}
\fonte{Elaborado pelos autores.}
\end{figure}

% ============================================================
\section{Análise do Protótipo do 1º CGEO}
\label{sec:prototipo}
% ============================================================

\subsection{Arquitetura e Contribuições}

A Divisão de Levantamento ``Gen Augusto Tasso Fragoso'', do 1º CGEO, desenvolveu o protótipo como prova de conceito do uso de LLMs locais na busca sobre o catálogo sob sua guarda \cite{cgeo2025prototipo}. O código é público, e os 22 casos de teste do protótipo formam a camada P do \textit{dataset} de avaliação (Seção~\ref{sec:camadas}). A Figura~\ref{fig:arquitetura_prototipo} mostra as quatro camadas.

\begin{figure}[h]
\centering
\begin{tikzpicture}[
  font=\footnotesize,
  node distance=0.65cm,
  layer/.style={rectangle, draw, rounded corners, minimum width=11cm, minimum height=1.0cm, align=center, fill=gray!6},
  arr/.style={-{Stealth[length=2.2mm]}, thick}
]
  \node[layer, fill=blue!8] (l1) {\textbf{Frontend} (React + TypeScript + Leaflet) \\ \scriptsize mapa interativo do Brasil};
  \node[layer, below=of l1] (l2) {\textbf{Backend} (Node.js + Express + TypeScript) \\ \scriptsize Pré-processamento (dicionário \textit{hardcoded} de termos)};
  \node[layer, below=of l2, fill=orange!10] (l3) {\textbf{LLM} (Ollama / Phi-4 14B) \\ \scriptsize Saída Estruturada via Instructor + Zod};
  \node[layer, below=of l3, fill=green!8] (l4) {\textbf{Banco de Dados} (PostgreSQL + PostGIS) \\ \scriptsize \textit{full-text} + busca espacial + GeoJSON};
  \draw[arr] (l1) to node[right, font=\scriptsize] {POST /api/natural-search} (l2);
  \draw[arr] (l2) to node[right, font=\scriptsize] {parâmetros de busca} (l3);
  \draw[arr] (l3) to node[right, font=\scriptsize] {SQL parametrizado} (l4);
  \draw[arr, dashed] (l4.east) to ++(0.6,0) |- node[right, font=\scriptsize, pos=0.25] {resultados (GeoJSON)} (l1.east);
\end{tikzpicture}
\caption{Arquitetura do protótipo de busca do 1º CGEO}
\label{fig:arquitetura_prototipo}
\fonte{Adaptado de \citeauthoronline{cgeo2025prototipo} (\citeyear{cgeo2025prototipo}).}
\end{figure}

O \textit{frontend}, em React com Leaflet, mostra os resultados num mapa do Brasil. O \textit{backend}, em Node.js com Express, recebe a consulta pela rota \texttt{POST /api/natural-search}. Antes do LLM, um pré-processador normaliza variações de grafia com um dicionário fixo no código (``25k'' $\rightarrow$ ``1:25.000''). O núcleo é o Phi-4 14B \cite{abdin2024phi4}, executado via Ollama, com Saída Estruturada pela biblioteca Instructor e \textit{schema} em Zod. O \textit{schema} pede um campo \texttt{reasoning}, em que o modelo justifica a extração, e os doze campos de busca (\texttt{keyword}, \texttt{scale}, \texttt{productType}, \texttt{state}, \texttt{city}, \texttt{supplyArea}, \texttt{project}, \texttt{publicationPeriod}, \texttt{creationPeriod}, \texttt{sortField}, \texttt{sortDirection} e \texttt{limit}). O \textit{prompt} traz exemplos resolvidos (\textit{few-shot prompting}). Se o LLM falhar mesmo após três novas tentativas, um \textit{fallback} por expressões regulares tenta extrair do texto parte dos campos (escala, tipo de produto, área de suprimento, projeto, ordenação, limite e períodos). O PostgreSQL/PostGIS combina busca textual com operações espaciais e devolve GeoJSON. O repositório traz o esquema do banco (\texttt{init.sql}, quatro tabelas), mas não os dados do acervo.

O protótipo demonstrou que (i) LLMs locais são tecnicamente viáveis no domínio cartográfico, (ii) o Phi-4 14B extrai parâmetros de busca relevantes de consultas em português, (iii) \textit{few-shot prompting} com Saída Estruturada produz resultados utilizáveis, o que este trabalho mediu na Seção~\ref{sec:abordagens}, e (iv) a integração do LLM com o PostgreSQL/PostGIS para busca espacial funciona em ambiente controlado.

\subsection{Limitações}
\label{sec:limitacoes}

O cliente apontou seis limitações do protótipo, que definem as frentes de aprimoramento:
\begin{enumerate}
  \item \textbf{\textit{Backend} em Node.js.} As bibliotecas de integração com LLMs (LangChain, LlamaIndex, Instructor) têm versões mais completas em Python.
  \item \textbf{Apenas Saída Estruturada.} O modelo devolve sempre a estrutura predefinida e não pode encadear ferramentas nem reagir ao resultado de operações anteriores.
  \item \textbf{Modelo único e defasado.} Com um só modelo, não há como comparar alternativas e escolher a mais adequada, num ecossistema que muda depressa.
  \item \textbf{Dicionário fixo no código.} O pré-processador aplica 159 pares de substituição antes de o texto chegar ao LLM. Variações não previstas no desenvolvimento (``dois cgeo'', ``25 mil'') falham em silêncio. A manutenção é manual, a partir dos erros vistos no \textit{log}. Um termo do dicionário em contexto ambíguo pode ainda deturpar uma consulta legítima.
  \item \textbf{\textit{Fallback} por expressões regulares.} Sua presença indica falhas do LLM que não foram documentadas de forma sistemática, o que dificulta a análise de causa.
  \item \textbf{Ausência de métricas de avaliação.} Os 22 casos de teste são pontuados por pesos atribuídos \textit{ad hoc} aos campos esperados, sem penalizar campos preenchidos indevidamente, e não há protocolo para comparar modelos, \textit{prompts} ou configurações.
\end{enumerate}

% ============================================================
\section{Requisitos}
\label{sec:requisitos}
\label{sec:requisitos-resumo}
% ============================================================

Os requisitos vieram da análise do protótipo e das conversas com o cliente. O Quadro~\ref{quadro:requisitos-resumo} traz cada um, a seção que o atende e como foi verificado. A latência (RNF2) depende do \textit{hardware} e da carga do ambiente de implantação, que estão fora do escopo. Por isso ela é reportada por ambiente (Seção~\ref{sec:ambiente}), sem limiar fixado \textit{a priori}.

\begin{quadro}[htbp!]
\centering
\caption{Requisitos da solução, onde são atendidos e como foram verificados}
\label{quadro:requisitos-resumo}
\footnotesize
\begin{tabular}{|c|>{\raggedright\arraybackslash}p{4.7cm}|>{\raggedright\arraybackslash}p{2.6cm}|>{\raggedright\arraybackslash}p{5.2cm}|}
\hline
\textbf{Requisito} & \textbf{Descrição} & \textbf{Atendido em} & \textbf{Verificação} \\
\hline
RF1 & Receber consulta em linguagem natural por rota REST & Seção~\ref{sec:api} & Demonstrado na API sobre os dados sintéticos; testes automatizados da API \\
\hline
RF2 & Traduzir a consulta nos doze parâmetros do \textit{schema} & Seção~\ref{sec:orquestracao} & Avaliado no Capítulo~\ref{cap:resultados} \\
\hline
RF3 & Executar a busca no PostgreSQL/PostGIS & Seção~\ref{sec:persistencia} & Demonstrado sobre os dados sintéticos; não avaliado \\
\hline
RF4 & Selecionar o modelo por requisição (só em desenvolvimento) & Seção~\ref{sec:api} & Implementado e coberto por teste da API; fora das métricas \\
\hline
RF5 & Registrar, por consulta, os parâmetros, as chamadas, os tempos e os erros, para o cálculo das métricas & Seção~\ref{sec:pipeline-avaliacao} & Registros e manifestos de todas as rodadas \\
\hline
RNF1 & Execução totalmente local & Seções~\ref{sec:orquestracao} e~\ref{sec:selecao-modelos} & Modelos locais; nuvem só como referência \\
\hline
RNF2 & Latência compatível com uso interativo & Seção~\ref{sec:ambiente} & Medida por ambiente, sem limiar fixado \\
\hline
RNF3 & Conformidade com a ET-PCDG e o Perfil MGB na estrutura dos metadados & Seção~\ref{sec:persistencia} & Esquema do protótipo reaproveitado; conformidade não verificada à parte \\
\hline
RNF4 & Reprodutibilidade dos experimentos & Seção~\ref{sec:reprodutibilidade} & Manifestos, \textit{hashes} e repontuação \\
\hline
RNF5 & Manutenibilidade & Seção~\ref{sec:harness} & Testes automatizados e integração contínua \\
\hline
RNF6 & Substituição do modelo por configuração & Seção~\ref{sec:harness} & Modelo e abordagem escolhidos por configuração \\
\hline
\end{tabular}%
\fonte{Elaborado pelos autores.}
\end{quadro}

% ============================================================
\section{Planejamento da Migração}
% ============================================================

A estratégia foi a \textbf{substituição dirigida}. Os componentes responsáveis pelas limitações da Seção~\ref{sec:limitacoes} foram substituídos, e o esquema do banco, os parâmetros de busca e a lógica das consultas espaciais foram aproveitados. O Quadro~\ref{quadro:mapa-limitacoes} liga cada limitação à solução adotada, e o Quadro~\ref{quadro:componentes} classifica cada componente do sistema final.

\begin{quadro}[h]
\centering
\caption{Mapeamento entre limitações do protótipo e soluções adotadas}
\label{quadro:mapa-limitacoes}
\small
\begin{tabular}{|c|>{\raggedright\arraybackslash}p{4.2cm}|>{\raggedright\arraybackslash}p{4.2cm}|>{\raggedright\arraybackslash}p{3.2cm}|}
\hline
\textbf{\#} & \textbf{Limitação} & \textbf{Solução adotada} & \textbf{Componente afetado} \\
\hline
1 & \textit{Backend} em Node.js & Reimplementação em Python com FastAPI & Nova API REST \\
\hline
2 & Apenas Saída Estruturada & \textit{Tool Calling} nativo via LangChain & Serviço de LLM \\
\hline
3 & Modelo único (Phi-4 14B) & Avaliação comparativa de quatro modelos atuais (Qwen 3 4B, Gemma 4 E2B, Gemma 4 E4B, Mistral Nemo 12B) & Configuração e testes \\
\hline
4 & Pré-processamento manual & Normalização delegada ao LLM, orientada pelas descrições do \textit{schema} da ferramenta & Remoção do pré-processador \\
\hline
5 & \textit{Fallback} via regex & Falhas registradas como métrica, sem \textit{fallback} oculto & Serviço de LLM \\
\hline
6 & Ausência de métricas & \textit{Dataset} auditado com gabarito + \textit{pipeline} de avaliação automatizado & Novo módulo de avaliação \\
\hline
\end{tabular}%
\fonte{Elaborado pelos autores.}
\nota{Soluções da v1. A v3 acrescenta, nas linhas 4 e 5, ferramentas auxiliares e o retorno de um validador (Seção~\ref{sec:orquestracao}).}
\end{quadro}

\begin{quadro}[h]
\centering
\caption{Classificação dos componentes do sistema final em relação ao protótipo}
\label{quadro:componentes}
\small
\begin{tabular}{|>{\raggedright\arraybackslash}p{4cm}|c|>{\raggedright\arraybackslash}p{6.5cm}|}
\hline
\textbf{Componente} & \textbf{Categoria} & \textbf{Observação} \\
\hline
Esquema do banco PostgreSQL/PostGIS & Reaproveitado & Mesmo esquema do protótipo (\texttt{init.sql}, quatro tabelas), sem alterações \\
\hline
Conjunto de parâmetros de busca & Reaproveitado & Os mesmos doze campos, compatíveis com a lógica das consultas SQL; enumerados recortados (9 tipos de produto e 9 projetos, contra 30 e 15 no protótipo) e grafados conforme a ET-PCDG \\
\hline
Lógica das consultas SQL & Reaproveitado & Portada de \texttt{search.ts}; direção de ordenação restrita a \texttt{ASC} ou \texttt{DESC}; aceita períodos com um só limite; sem o filtro por retângulo do mapa e a paginação, próprios do \textit{frontend} \\
\hline
\textit{Frontend} React/Leaflet & Fora do escopo & Não é parte deste trabalho \\
\hline
API REST & Reescrito & Node.js/Express $\to$ Python/FastAPI \\
\hline
Serviço de LLM & Reescrito & Saída Estruturada (Instructor/Zod) $\to$ \textit{Tool Calling} (LangChain) \\
\hline
Camada de acesso ao banco & Reescrito & \texttt{pg-promise} $\to$ \texttt{psycopg} \\
\hline
Pré-processamento de termos & Eliminado & Função transferida ao LLM \\
\hline
\textit{Fallback} com regex & Eliminado & Falhas viram métricas \\
\hline
Dados sintéticos do banco & Novo & Apenas para a demonstração ponta a ponta (Seção~\ref{sec:persistencia}) \\
\hline
\textit{Dataset} de avaliação & Novo & 310 consultas com gabarito auditado (Cap.~\ref{cap:metodologia} e Ap.~\ref{apendice:dataset}) \\
\hline
\textit{Pipeline} de avaliação & Novo & Execução, pontuação, conferência e relatórios (Seção~\ref{sec:pipeline-avaliacao}) \\
\hline
Ferramentas auxiliares e validador (v3) & Novo & No pacote, no avaliador e na API, como abordagem padrão (Seção~\ref{sec:v3}) \\
\hline
Registro de abordagens, testes e integração contínua & Novo & Registro único das configurações avaliadas, testes de regressão e guia de extensão (Seção~\ref{sec:harness}) \\
\hline
\end{tabular}%
\fonte{Elaborado pelos autores.}
\end{quadro}

% ============================================================
\section{Arquitetura da Solução v1}
\label{sec:desenvolvimento}
% ============================================================

A Figura~\ref{fig:arquitetura_proposta} mostra as camadas da solução. A v2 e a v3 mudam só a camada de tradução (Seções~\ref{sec:v2} e~\ref{sec:v3}). O padrão da API é hoje o \textit{Tool Calling} v3 com o Gemma 4 E4B (Seção~\ref{sec:harness}).

\begin{figure}[h]
\centering
\begin{tikzpicture}[
  font=\footnotesize,
  node distance=0.6cm,
  layer/.style={rectangle, draw, rounded corners, minimum width=11.5cm, minimum height=1.0cm, align=center, fill=gray!6},
  half/.style={rectangle, draw, rounded corners, minimum width=5.5cm, minimum height=1.2cm, align=center, fill=gray!6},
  arr/.style={-{Stealth[length=2.2mm]}, thick},
  darr/.style={{Stealth[length=2.2mm]}-{Stealth[length=2.2mm]}, thick}
]
  \node[layer, fill=blue!8] (l1) {\textbf{Cliente} (aplicação consumidora da API, HTTP/JSON)};
  \node[layer, below=of l1] (l2) {\textbf{FastAPI} (Python) \\ \scriptsize \texttt{POST /api/search} e \texttt{GET /api/health}; validação Pydantic};
  \node[layer, below=of l2, fill=orange!10] (l3) {\textbf{\textit{Pipeline} + LangChain} (\texttt{ChatOllama.bind\_tools()}) \\ \scriptsize tradução por \textit{Tool Calling}; execução da ferramenta; resposta final \\ \scriptsize padrão da API é a v3, com ferramentas auxiliares e validador num laço de chamadas (Seção~\ref{sec:v3})};
  \node[half, below=1.1cm of l3.south west, anchor=north west, fill=purple!8] (l4) {\textbf{Ollama} (local) \\ \scriptsize Qwen 3 4B \;|\; Gemma 4 E2B \;|\; Gemma 4 E4B \\ \scriptsize Mistral Nemo 12B};
  \node[half, below=1.1cm of l3.south east, anchor=north east, fill=green!8] (l5) {\textbf{PostgreSQL + PostGIS} \\ \scriptsize \texttt{buscar\_catalogo} com SQL parametrizada \\ \scriptsize saída em GeoJSON};
  \draw[arr] (l1) to node[right, font=\scriptsize] {\texttt{POST /api/search} (JSON)} (l2);
  \draw[arr] (l2) to node[right, font=\scriptsize] {consulta} (l3);
  \draw[darr] (l4.north) to node[right, font=\scriptsize] {\textit{tool call} (JSON)} (l4.north |- l3.south);
  \draw[darr] (l5.north) to node[right, font=\scriptsize] {parâmetros / produtos} (l5.north |- l3.south);
  \draw[arr, dashed] (l2.east) to ++(0.6,0) |- node[right, font=\scriptsize, pos=0.25] {resposta JSON} (l1.east);
\end{tikzpicture}
\caption{Arquitetura da solução aprimorada de busca}
\label{fig:arquitetura_proposta}
\fonte{Elaborado pelos autores.}
\end{figure}

\subsection{API}
\label{sec:api}

A API REST usa o FastAPI~\cite{fastapi2018}, pela validação das requisições com Pydantic e pela documentação gerada automaticamente. A rota \texttt{POST /api/search} recebe a consulta e devolve os parâmetros extraídos, os produtos em GeoJSON, a resposta ao usuário, as latências e os erros. A rota \texttt{GET /api/health} informa a versão do Ollama, o estado do banco, o modelo e a abordagem no ar. O \textit{pipeline} de avaliação chama diretamente a camada de tradução, sem passar pela API, para que a latência medida seja a da chamada ao modelo.

\paragraph{Seleção de modelo.} A requisição pode escolher o modelo (RF4) só no ambiente de desenvolvimento. Em operação, o modelo é fixado por configuração, porque trocar de modelo entre requisições recarrega pesos na VRAM.

\paragraph{Tratamento de falhas.} Na v1 e na v2, cada falha do LLM é registrada como métrica, sem novas tentativas nem \textit{fallback} (limitação 5). As falhas se dividem em erros de infraestrutura (tempo esgotado, servidor ou modelo indisponível) e erros do modelo (argumentos fora do \textit{schema}, chamada a ferramenta inexistente). Um \textit{fallback} de produção fica como trabalho futuro (Seção~\ref{sec:trabalhos-futuros}).

\subsection{Orquestração e Ferramenta de Busca}
\label{sec:orquestracao}

A orquestração usa o \texttt{ChatOllama} do LangChain~\cite{langchain2022}, com a configuração da Seção~\ref{sec:config-modelos}. Na v1, a tradução é uma única chamada ao modelo, e o ciclo completo (executar a ferramenta e devolver o resultado ao modelo para a resposta final) é escrito no próprio módulo de \textit{pipeline}.

A ferramenta \texttt{buscar\_catalogo} tem os doze parâmetros de busca, cada um com tipo e, quando cabe, a lista fechada de valores válidos (o enumerado). Dez deles trazem na descrição instruções de normalização. O \textit{schema} fica num único módulo, que também serve para validar os argumentos do modelo e os enumerados do \textit{dataset}, e o seu \textit{hash} vai para o manifesto de cada rodada. Sem dicionário nem exemplos resolvidos no \textit{prompt}, o modelo normaliza os termos guiado só pelas descrições, que substituem as 159 entradas do dicionário do protótipo (limitação 4).

\begin{sloppypar}
A v1 expõe só \texttt{buscar\_catalogo}, e a v2 acrescenta \texttt{recusar\_consulta} (Seção~\ref{sec:v2}). A v3 (Seção~\ref{sec:v3}) acrescenta \texttt{identificar\_\allowbreak nome}, \texttt{normalizar\_\allowbreak codigo}, \texttt{normalizar\_\allowbreak escala}, \texttt{resolver\_\allowbreak periodo} e \texttt{pedir\_\allowbreak esclarecimento}. A v3 devolve parte da normalização ao código, em ferramentas que o modelo consulta quando precisa e num validador que confere a saída, e a consulta chega ao modelo sem reescrita prévia.
\end{sloppypar}

\subsection{Persistência}
\label{sec:persistencia}

O banco usa, sem alterações, o esquema do protótipo, em PostgreSQL com PostGIS~\cite{postgis2001}. Como o protótipo não traz os dados do acervo, a demonstração ponta a ponta usa dados sintéticos, gerados por \textit{script} com semente fixa (27 unidades da federação representadas por retângulos numa grade, 14 municípios, cinco áreas de suprimento e 117 produtos). \textbf{Nenhuma métrica deste trabalho depende do banco}, porque a avaliação incide sobre os parâmetros extraídos pelo modelo. O Docker Compose sobe o PostgreSQL 16 com PostGIS 3.4.

A lógica portada do \texttt{search.ts} combina busca textual, códigos MI e INOM, filtros, períodos, interseções espaciais, ordenação e limite, com os ajustes do Quadro~\ref{quadro:componentes}. As consultas são parametrizadas, e o modelo nunca escreve SQL.

Os enumerados do novo \textit{schema} grafam alguns valores de modo diferente do protótipo (acentos, travessões e nomes). Os dados sintéticos seguem a grafia nova, e ligar a API ao banco do acervo exige conferir a grafia guardada nele. A resolução dos campos geográficos em geometrias fica fora do escopo (Seção~\ref{sec:exclusoes-avaliacao}).

\subsection{Consultas sem Resposta no Acervo}
\label{sec:sem-resposta}

O sistema distingue três desfechos, e em nenhum deles inventa um produto:
\begin{enumerate}
  \item \textbf{Consulta fora do domínio do catálogo} (``cartas de Marte''). Não há busca. Na v1, o modelo não chama a ferramenta e responde em uma frase por que o pedido está fora do escopo. No \textit{Tool Calling} v2 e v3, ele chama \texttt{recusar\_consulta}, cujo argumento \texttt{motivo} é a frase dirigida ao usuário (na Saída Estruturada v2 e v3, a recusa é o campo \texttt{fora\_do\_escopo}). Na v3, a primeira recusa de uma consulta que cita um critério do catálogo volta ao modelo para confirmação (Seção~\ref{sec:v3}). A categoria F do \textit{dataset} mede esse comportamento.
  \item \textbf{Busca com produtos.} A lista vem sempre do banco, e o modelo só a resume.
  \item \textbf{Busca legítima sem produto.} A API devolve a lista vazia e uma mensagem fixa com os critérios usados, que sugere retirar ou ampliar algum deles, sem nova chamada ao modelo.
\end{enumerate}

Nos pedidos com valores que a ferramenta não representa (``escala 1:10.000.000'', ``cartas do estado 42''), o texto não determina o comportamento correto. O \textit{dataset} os trata como casos observacionais, fora das métricas principais (Seção~\ref{sec:observacionais}).

% ============================================================
\section{\textit{Pipeline} de Avaliação}
\label{sec:pipeline-avaliacao}
% ============================================================

O \textit{pipeline} de avaliação produz os números do Capítulo~\ref{cap:resultados}. Ele gera, da mesma fonte, o arquivo de consultas avaliado e as listagens do Apêndice~\ref{apendice:dataset}, validados contra o \textit{schema}, e roda as checagens automáticas da auditoria do gabarito (Seção~\ref{sec:auditoria}). Executa as consultas em ordem fixa, de forma sequencial e retomável, guarda a resposta, os tempos e os erros de cada consulta e repetição e grava um manifesto de reprodutibilidade por rodada (RNF4). A pontuação é sempre recalculada contra o gabarito vigente, e uma correção do gabarito não exige nova execução dos modelos. Antes da consolidação, ele confere os \textit{hashes} e a completude das rodadas. Depois gera as tabelas, as figuras e os testes estatísticos, inclusive os dos lotes e do teste A/B. A referência em nuvem roda no mesmo avaliador, pelo Groq (Apêndice~\ref{apx:specs-groq}). Os comandos estão no Apêndice~\ref{apx:comandos}.

\subsection{Registro de Abordagens e Implantação}
\label{sec:harness}

Um módulo único, o registro de abordagens, reúne as nove configurações avaliadas no estudo, da solução v1 às especificações v2 e v3 e ao método do protótipo (Quadro~\ref{quadro:desenho}). Os módulos das versões avaliadas ficam congelados, e um teste de regressão confere que o código atual reproduz os \textit{hashes} de \textit{prompt} e de ferramenta de todas as rodadas válidas.

A API escolhe a abordagem e o modelo por configuração (RNF6). O padrão é o \textit{Tool Calling} v3 com o Gemma 4 E4B (Seção~\ref{sec:recomendacao}), e a v1 continua disponível. O \texttt{/api/health} mostra os \textit{hashes} de \textit{prompt} e de ferramenta e as rodadas feitas com a mesma configuração, o que permite verificar se o que está no ar foi o avaliado. Na API, o pedido de esclarecimento da v3 continua respondido pelo usuário simulado da avaliação (Seção~\ref{sec:v3}), e levá-lo ao usuário real exige conversa em vários turnos (Seção~\ref{sec:trabalhos-futuros}). O índice de folhas do BDGEx, de que as ferramentas da v3 dependem, não está no repositório. Sem ele, a API sobe em modo degradado, com aviso, e o teste de regressão pula os casos da v3 (Apêndice~\ref{apendice:codigo}).

Um \textit{Dockerfile} empacota a API. Os testes automatizados cobrem o \textit{schema}, as consultas SQL, as métricas e a API, sem depender do Ollama nem do banco (RNF5), e a integração contínua os roda a cada mudança (Apêndice~\ref{apx:comandos}). Um guia de extensão exige que cada versão nova seja testada num conjunto ainda não usado, e uma v4 precisaria de um terceiro lote. Os autores implementaram muitas partes do sistema, e o restante do código e dos testes foi escrito com apoio do Claude Code, com os modelos Claude Opus 4.6 e Claude Opus 5.5, a partir da especificação deles (Tabela~\ref{tab:uso-ia}).

% ============================================================
\section{Cronograma}
\label{sec:cronograma}
% ============================================================

O projeto foi executado entre abril e outubro de 2026, em oito fases (Quadro~\ref{quadro:cronograma}). As fases F0, F1 e F5 correspondem aos pacotes de construção da EAP, e as fases F2, F3, F4 e F6 à avaliação, detalhada no Capítulo~\ref{cap:metodologia}. A F6 aparece dividida pelas datas do histórico do repositório.

\begin{quadro}[h]
\centering
\caption{Cronograma de execução do projeto}
\label{quadro:cronograma}
\small
\begin{tabular}{|c|p{6.6cm}|c|p{3.4cm}|}
\hline
\textbf{Fase} & \textbf{Atividade} & \textbf{Período} & \textbf{Marco} \\
\hline
F0 & Análise do protótipo: código-fonte, dependências, esquema do banco e limitações documentadas & abr/2026 & Limitações mapeadas \\
\hline
F1 & Elicitação de requisitos e planejamento da migração (limitação $\to$ solução; componentes reaproveitados, reescritos e novos) & mai/2026 & Requisitos e plano de migração \\
\hline
F2 & Referencial teórico; seleção dos modelos; definição das métricas e do protocolo de avaliação & mai--jun/2026 & Protocolo de avaliação \\
\hline
F3 & Construção do \textit{dataset} de avaliação: camadas P, N e G e respectivo gabarito (Apêndice~\ref{apendice:dataset}) & jun--jul/2026 & 310 consultas anotadas \\
\hline
F4 & Prova de conceito do \textit{Tool Calling} sobre consultas representativas, com modelo local e em nuvem & ago/2026 & Viabilidade do \textit{Tool Calling} demonstrada \\
\hline
F5 & Implementação: API, orquestração, ferramenta \texttt{buscar\_catalogo}, persistência e \textit{pipeline} de avaliação, com testes automatizados & set/2026 & Sistema ponta a ponta \\
\hline
F6a & Rodadas completas das 310 consultas, validação na estação de referência, referência em nuvem e auditoria do \textit{dataset} & 14--26/set/2026 & Resultados da v1 \\
\hline
F6b & Linhas de base com Saída Estruturada e método do protótipo; lote 1 e especificação v2 & 28/set/2026 & Comparação das abordagens \\
\hline
F6c & Rodadas do lote 1; v3 congelada; lote 2 e teste da v3 & 29/set/2026 & Teste da v3 \\
\hline
F6d & Extensão da v3 a outros modelos; registro de abordagens e API com a v3; teste A/B; latência na estação & 29--30/set/2026 & Recomendação \\
\hline
F7 & Redação dos resultados e da conclusão; preparação da defesa & set--out/2026 & Defesa (out/2026) \\
\hline
\end{tabular}%
\fonte{Elaborado pelos autores.}
\end{quadro}
